\documentclass{scrartcl}
\usepackage{graphicx}
\usepackage{float}
\usepackage{amsmath}
\usepackage{verbatim}
\usepackage{caption}
\usepackage{subcaption}

\newcommand{\figwidth}{4in}
\newcommand\tab[1][0.5cm]{\hspace*{#1}}
\renewcommand\thesubsection{\thesection.\alph{subsection}}

\title{5SWC0 - Linear and Planar Motors for High-precision Systems}
\subtitle{Lab Assignment Group 1}

\begin{document}

\begin{titlepage}
\author{J.M. van Dam (ID: 0806678)\\j.m.v.dam@student.tue.nl\\Department of Mechanical Engineering
\and \\
M. Kleijer (ID:0813485)\\m.kleijer@student.tue.nl\\Department of Electrical Engineering}


\clearpage\maketitle
\thispagestyle{empty}
\end{titlepage}

\newpage
\setcounter{page}{1}

\section{Introduction}
	\label{1}
	\input{Chapters/Introduction.tex}

\section{EMF}
\tab In the first experiment, the set-up has no external excitation, except for the manual pushing of the actuator. This means that the voltage measured with the dSpace systems equals the induced EMF.

	\subsection{Question 1}
    \label{2a}
    \input{Chapters/Question_1.tex}
    
    \subsection{Question 2}
    \label{2b}
    \input{Chapters/Question_2.tex}
    
    \subsection{Question 3}
    \label{2c}
    \input{Chapters/Question_3.tex}
    
    \subsection{Question 4}
    \label{2d}
    \input{Chapters/Question_4.tex}

\section{Phase resistances}
\tab In this experiment the phase resistances of the machine set are measured with the LCR meter.

	\subsection{Question 5}
    \label{3a}
    \input{Chapters/Question_5.tex}

\section{Phase inductances}
\tab Next, the LCR meter is used to measure the phase inductances.

	\subsection{Question 6}
    \label{4a}
    \input{Chapters/Question_6.tex}
    
    \subsection{Question 7}
    \label{4b} 
    \input{Chapters/Question_7.tex}

    \subsection{Question 8}
    \label{4c}
    \input{Chapters/Question_8.tex}

\section{Friction force}
\tab For this experiment the machine is actuated with certain motion profiles from which the static friction or Coulomb friction and dynamic friction or damping can be derived. The motor constant calculated via the EMF measurements is implemented into the dSpace measurement environment.

    \subsection{Question 9}
    \label{5a}
    \input{Chapters/Question_9.tex}
    
    \subsection{Question 10}
    \label{5b}
    \input{Chapters/Question_10.tex}
    
    \subsection{Question 11}
    \label{5c}
    \input{Chapters/Question_11.tex}
    
\section{Cogging force}
\tab Besides the friction force, the cogging force also acts as a disturbance force. The latter has yet to be mapped, which is done in this experiment.

	\subsection{Question 12}
    \label{6a}
    \input{Chapters/Question_12.tex}
    
	\subsection{Question 13}
    \label{6b}
    \input{Chapters/Question_13.tex}
    
	\subsection{Question 14}
    \label{6c}
    \input{Chapters/Question_14.tex}
    
\section{Commutation angle}
\tab In this experiment the effect of a phase error with respect to the optimum commutation angle on the force production of the linear motor is investigated; it is known that a perpendicular commutation angle is most efficient.

	\subsection{Question 15}
    \label{7a}
    \input{Chapters/Question_15.tex}
    
    \subsection{Question 16}
    \label{7b}
    \input{Chapters/Question_16.tex}
    
    \subsection{Question 17}
    \label{7c}
    \input{Chapters/Question_17.tex}
    
\section{Feed-forward controller}
\tab Finally, the filtered waveform for the cogging force is applied to the look-up table for the provided feed-forward controller. The effect of such a controller is investigated during this experiment.

	\subsection{Question 18}
    \label{8a}
    \input{Chapters/Question_18.tex}
    
\end{document}
